\section{Thiết lập thí nghiệm}
\label{sec:thiet-lap}

\subsection{Dữ liệu và thuật toán}

Mỗi tác vụ dùng một bộ dữ liệu công khai đủ lớn (Bảng~\ref{tab:datasets}). HIGGS được chọn vì có bản đầy đủ 11 triệu dòng, dùng cho phần scalability; Covertype là bộ duy nhất trùng với bài gốc. Bản ``sạch'' của Beijing bỏ 9,2\% số dòng có ô thiếu tự nhiên, để baseline có Completeness bằng 1 và ô thiếu tự nhiên không lẫn với ô thiếu do polluter tạo ra.

\begin{table}[htbp]
\centering
\small
\setlength{\tabcolsep}{3.5pt}
\begin{tabular}{@{}lllrll@{}}
\toprule
Tác vụ & Bộ dữ liệu & Mục tiêu & Dòng $\times$ cột & Thuật toán & Thước đo \\
\midrule
Regression     & Beijing Air Quality \citep{zhang2017cautionary} & PM2.5 & 382.168 $\times$ 42   & \makecell[l]{Ridge,\\Gradient Boosting} & R$^2$ \\
Classification & HIGGS mẫu \citep{baldi2014searching}            & 2 lớp & 1.000.000 $\times$ 29 & \makecell[l]{Logistic Regression,\\Random Forest} & F1-macro \\
Clustering     & Covertype \citep{blackard1999comparative}       & 7 lớp & 581.012 $\times$ 55   & \makecell[l]{k-Means,\\Gaussian Mixture} & AMI \\
\bottomrule
\end{tabular}
\caption{Dữ liệu và thuật toán theo tác vụ. Mỗi tác vụ có một thuật toán đơn giản (tuyến tính hoặc dựa trên tâm cụm) và một thuật toán phức tạp hơn.}
\label{tab:datasets}
\end{table}

Mỗi tác vụ dùng hai thuật toán khác nhau để thấy ảnh hưởng của thuật toán ngay trong cùng tác vụ; tất cả được cài đặt bằng scikit-learn \citep{pedregosa2011scikit}. Trước khi huấn luyện, ô thiếu được điền bằng trung vị của tập train và các đặc trưng được chuẩn hóa z-score. Classification và regression chia train/test theo tỷ lệ 80/20. Clustering dùng số cụm bằng số lớp thật và so các cụm với nhãn bằng AMI. Hai cấu hình được chọn dựa trên số đo thay vì dùng mặc định: Random Forest học mỗi cây trên 10\% số dòng (F1 chỉ thấp hơn 0,008 mà nhanh gấp 6,6 lần), và Gaussian Mixture dùng hiệp phương sai đường chéo với 10 lần khởi tạo (độ lệch chuẩn của AMI theo seed giảm từ 0,045 xuống 0,010). Các siêu tham số còn lại dùng giá trị mặc định của scikit-learn, trừ \texttt{max\_iter=1000} (Logistic Regression) và \texttt{n\_init=10} (k-Means).

\subsection{Ma trận thí nghiệm}

Mỗi tác vụ chạy mọi tổ hợp của (chiều chất lượng áp dụng được) $\times$ (2 thuật toán) $\times$ (4 mức $\lambda$) $\times$ (3 seed: 42, 43, 44). Regression không có Class Balance và clustering không có Target Accuracy, nên tổng số lượt là $4{\cdot}2{\cdot}4{\cdot}3 + 5{\cdot}2{\cdot}4{\cdot}3 + 4{\cdot}2{\cdot}4{\cdot}3 = 96 + 120 + 96 = 312$ (lần lượt cho regression, classification và clustering). Toàn bộ 312 lượt chạy trên Google Colab (CPU 2 nhân, 13,6 GB RAM) trong 414,5 phút.

Trước khi chạy, các polluter được kiểm tra bằng cách làm bẩn một mẫu 100.000 dòng rồi đo lại bằng chính profiler (Hình~\ref{fig:h1}). Điểm đo được bám sát đường lý thuyết: $1-\lambda$ cho Completeness, Uniqueness và Target Accuracy của classification; $1-\sqrt{\lambda}\sqrt{2/\pi}$ cho accuracy của cột số. Feature Accuracy của Beijing và Covertype nằm trên đường lý thuyết ở các mức $\lambda$ thấp và trung bình, vì hai bộ này có nhóm cột one-hot, nơi điểm giảm theo $1-\lambda$ chứ không theo công thức nhiễu Gaussian. Ngoại lệ đáng chú ý là Class Balance của Covertype: điểm \emph{tăng} khi dữ liệu bị làm bẩn (thảo luận ở Mục~\ref{sec:kb-balance}).

\begin{figure}[htbp]
\centering
\hinh[0.9\linewidth]{h1_pollution_check.png}
\caption{Điểm profile đo lại sau khi làm bẩn, theo $\lambda$ (đường liền) so với lý thuyết (đường đứt). Hàng: tác vụ; cột: chiều bị làm bẩn.}
\label{fig:h1}
\end{figure}

Hiệu năng trên dữ liệu sạch (Hình~\ref{fig:h2}) là mốc để tính $\Delta$. Bốn thuật toán của regression và classification gần như không đổi theo seed. Hai thuật toán clustering dao động theo seed, nhiều nhất là k-Means (AMI $0{,}172 \pm 0{,}017$). Đây là lý do $\Delta$ được tính so với baseline có cùng seed.

\begin{figure}[htbp]
\centering
\hinh[0.65\linewidth]{h2_baseline.png}
\caption{Hiệu năng baseline trên dữ liệu sạch (trung bình $\pm$ độ lệch chuẩn qua 3 seed).}
\label{fig:h2}
\end{figure}

\subsection{Cách đánh giá meta-model}
\label{sec:danh-gia}

\paragraph{Chia train/test theo nhóm.} Ba seed của cùng một cấu hình có profile gần như trùng nhau và $\Delta$ lệch nhau rất ít. Hai thuật toán cùng mức bẩn còn có chung profile. Nếu chia ngẫu nhiên theo dòng thì mô hình có thể ``nhớ'' đáp án từ các dòng gần giống nhau. Vì vậy đề tài dùng \texttt{StratifiedGroupKFold} 5 fold với nhóm là (tác vụ, chiều, $\lambda$); các dòng $\lambda = 0$ của mỗi tác vụ gộp chung một nhóm, tổng cộng có 42 nhóm. Cách chia được lặp lại với 3 seed và mọi phương pháp dùng đúng cùng các fold. Nhờ vậy mỗi dòng có một dự đoán out-of-fold, tức là dự đoán khi mô hình chưa từng thấy mức bẩn đó.

\paragraph{Chỉ số.} MAE đo sai số tuyệt đối; Spearman đo khả năng xếp hạng đúng mức suy giảm (bộ dữ liệu nào bị suy giảm hiệu năng nhiều hơn); R$^2$ được báo kèm. Mọi chỉ số được tính riêng cho từng lần lặp rồi báo trung bình $\pm$ độ lệch chuẩn, tách theo tác vụ.

\paragraph{Ý nghĩa thống kê.} Độ lệch chuẩn qua 3 lần lặp chỉ phản ánh dao động do cách chia fold. Để biết meta-model hơn baseline có thật hay không, đề tài dùng \emph{bootstrap theo nhóm}: lấy mẫu lại có hoàn lại 42 nhóm (giữ nguyên số nhóm của mỗi tác vụ), lặp 1.000 lần, rồi lấy khoảng tin cậy 95\% của hiệu MAE. Đơn vị lấy mẫu là nhóm chứ không phải dòng, vì các dòng trong một nhóm phụ thuộc lẫn nhau.

\paragraph{Bài toán Good/Bad.} Để so sánh được với rule-based, mỗi lượt được gán nhãn thật ``Bad'' nếu hiệu năng mất hơn 10\% ($\Delta$ tương đối $> 0{,}1$). Dự đoán của ba phương pháp còn lại được chuyển thành Good/Bad theo cùng ngưỡng. Chỉ số chính là F1 của lớp Bad và balanced accuracy, kèm phân tích độ nhạy với ngưỡng 5\% và 20\%.
